\chapter*{Glossaire}
\label{chap:Glossaire}
\addstarredchapter{Glossaire}
\markboth{Glossaire}{}
\begin{itemize}
\item[•] \textbf{$m$-vec :} vecteur de $m$ éléments de $\mathbb{F}_{16}$, empaqueté à deux éléments par octet.

\item[•] \textbf{Limbe :} mot de 64 bits d'un accumulateur.

\item[•] \textbf{Voie (\emph{lane}) :} moitié de 128 bits d'un registre ymm.

\item[•] \textbf{Table de multiplication (multab) :} table de 16 octets donnant, pour une paire de nombres, les deux produits par un quartet.

\item[•] \textbf{Corps fini :} ensemble fini de nombres où l'on peut additionner, soustraire, multiplier et diviser (par exemple $\mathbb{F}_{16}$).\item \textbf{Forme quadratique :} expression de degré 2 en plusieurs variables, par exemple $x_1x_2+x_3^2$.\item \textbf{Empreinte (hachage) :} courte suite de nombres calculée à partir d'un message ; deux messages différents ont des empreintes différentes.\item \textbf{Sel :} valeur aléatoire ajoutée au message avant le calcul de l'empreinte.\item \textbf{Compilateur, assembleur :} le compilateur transforme le code C en assembleur, la suite d'instructions que le processeur exécute.\item \textbf{Passe (du compilateur) :} une étape d'optimisation du code (par exemple SRA, PRE, DSE).\item \textbf{Bootstrap :} méthode qui rejoue le calcul sur des tirages au hasard des mesures pour estimer son incertitude.\item \textbf{Fin de vecteur :} fin d'un $m$-vec qui ne remplit pas un registre complet de 32 octets.

\item[•] \textbf{Recouvrement :} chargement reculé du dernier registre, qui retraite des octets déjà traités.

\item[•] \textbf{SAFE :} variante d'A2 qui ne lit rien au-delà de l'entrée de $P_1$.

\item[•] \textbf{Tick TSC :} unité du compteur d'horodatage (\code{rdtscp}) ; ce n'est pas un cycle cœur.

\item[•] \textbf{Processus :} unité statistique : une exécution du banc, 3\,001 paires.

\item[•] \textbf{Doublon :} écritures de 32 octets supplémentaires par accumulateur dans le code généré, selon la construction.

\item[•] \textbf{E1, E1b :} analyses statiques du code généré (E1b : origine du doublon, exécutée le 25 octobre 2026 sous GCC~16.2.0).

\item[•] \textbf{Fait, inférence, hypothèse :} statuts des énoncés, notés \F{}, \I{}, \Hy{}.

\end{itemize}
\clearpage
